%-------------------------
% Rover Resume - Star Template
% Link: https://github.com/subidit/rover-resume
%
% CUSTOM COMMANDS \aux{} and \rside{} for part formatting
% eg. \subsection{Company Name \aux{Designation} \rside{Duration}}
% Double lined titles with \subsection{} and \subsubsection{}
% \linkicon gives external link icon
%-------------------------

\documentclass[11pt]{article}
\usepackage[T1]{fontenc}
\usepackage[utf8]{inputenc}

\usepackage[sfdefault,semibold,lf]{FiraSans}
\usepackage[a4paper,margin=1in]{geometry}

\usepackage{xcolor}
\definecolor{accent}{HTML}{141E61}

\newcommand{\linkicon}{
  \color{gray}{\footnotesize\raisebox{1pt}{\faIcon{external-link-alt}}}
}

\setcounter{secnumdepth}{0}
\pdfgentounicode=1


%====== LIST FORMATTING ======
\usepackage{enumitem}
\setlist[itemize]{nosep, left=0pt .. 1.5em}
\setlist[enumerate]{left=0pt..1.5em}
\setlist[description]{itemsep=0pt}


%====== TITLE FORMATTING ======
\usepackage{xhfill}
\usepackage{soul}
\newcommand{\midrule}[1]{
  \leavevmode
  \xrfill[.5ex]{1pt}[accent]~\so{#1}~\xrfill[.5ex]{1pt}[accent]
}

\usepackage{titlesec}
\titlespacing{\section}{0pt}{*4}{*1}
\titlespacing{\subsection}{0pt}{*3}{*0}
\titlespacing{\subsubsection}{0pt}{*0}{*0}
\titleformat{\section}{\large\bfseries\uppercase}{}{}{\midrule}
\titleformat*{\subsection}{\large\bfseries\scshape}


%====== CUSTOM COMMANDS ======
\newcommand{\aux}[1]{
  $|$ \space {\normalfont\itshape #1}%
}

\newcommand{\rside}[1]{
  \hfill {\normalfont\color{accent} #1}%
}


%====== FOOTER ======
\usepackage{lastpage}
\usepackage{fancyhdr}
\pagestyle{fancy}
\fancyhf{}
\renewcommand{\headrulewidth}{0pt}

% Replace 'Star Rover' with your name in next line
\fancyfoot[C]{\small\color{gray} Autorizzo l'uso dei miei dati personali in conformità al GDPR 679/16 - "Regolamento europeo sulla protezione dei dati personali". \\
\vspace{0.25cm} % Puoi regolare lo spazio
Emanuele Seminara -- Page \thepage\ of \pageref*{LastPage}}

% \pagestyle{empty} 
% uncomment previous line to hide page numbering


%====== PACKAGES ======
\usepackage{fontawesome5}
\usepackage{multicol}
\usepackage[bookmarks=false,hidelinks]{hyperref}
\usepackage[protrusion=true,expansion=true]{microtype}


%====== FINISH PREAMBLE ======

%====== DOCUMENT STARTS HERE ======

\begin{document}

%====== BANNER ======

\begin{center}
  {
    \fontsize{36}{12} 
    \fontseries{heavy}\selectfont
    \color{accent} 
    Emanuele Seminara % YOUR NAME HERE
  } \\ \medskip

  % {\color{gray}{\faHome}} 1 infinite loop, CA \quad

  \href{tel:+393204854761}{{\color{gray}{\faPhone}} +39 320 485 4761} \quad
  \href{mailto:seminara.emanuele96@gmail.com}{{\color{gray}{\faEnvelope}} seminara.emanuele96@gmail.com} \quad 
  \href{https://github.com/EmanueleSeminara}{{\color{gray}{\faGithub}} EmanueleSeminara} \quad \quad \quad \quad
  {\color{gray}{\faHome}} Palermo, Italia \quad
  \href{https://www.linkedin.com/in/emanuele-seminara/}{{\color{gray}{\faLinkedin}} emanuele-seminara} \quad
  \href{https://emanueleseminara.it/}{{\color{gray}{\faLink}} emanueleseminara.it} 
\end{center}

%====== BANNER END ======


\section{Chi sono}
%==========
\noindent \textbf{Ingegnere Informatico} con oltre \textbf{4 anni di esperienza professionale nell'IT}, di cui \textbf{3 nello sviluppo software} come \textbf{Full Stack Developer}, dal backend \textbf{Java/Spring} alle pipeline dati in \textbf{Python}. Esperienza su architetture a microservizi, piattaforme dati (\textbf{MongoDB}, \textbf{Kafka}) e pratiche \textbf{DevOps} su \textbf{Azure}, dalla gestione degli ambienti alle pipeline CI/CD. Background in \textbf{cybersecurity} e ideatore del progetto vincitore dell'hackathon interno aziendale, un portale HR basato su \textbf{AI}.

\section{Esperienze}
%==========
\subsection{Elca Informatique SA \rside{Palermo}}
\subsubsection{Full Stack Developer \aux{Associate Engineer $\rightarrow$ Engineer} \rside{Apr 2024 -- Ora}}
\begin{itemize}
  \item \textbf{Piattaforma di Data Hub (Progetto Attuale):} Contributo allo sviluppo di un data hub che centralizza e ridistribuisce i dati tra i diversi sistemi aziendali di un cliente assicurativo. Sviluppo di pipeline \textbf{ETL} in \textbf{Python} e \textbf{Java} su un data layer basato su \textbf{MongoDB} con streaming \textbf{Kafka}, su infrastruttura \textbf{Azure}.
  \item \textbf{Migrazione Sistema Gestionale:} Automazione della sincronizzazione dati tra il gestionale storico (progetti e ore) e il nuovo sistema \textbf{ERP} aziendale, nell'ambito di un processo di migrazione graduale, tramite script \textbf{Python} personalizzati.
  \item \textbf{Hackathon Interno per l'Innovazione HR (Vincitore):} Ideazione del progetto, selezione del team e guida delle decisioni chiave su un portale HR basato su \textbf{AI} per screening dei CV, profilazione dei candidati e generazione automatica di test tecnici. Sviluppo del backend in \textbf{Python (FastAPI)}, delle REST API e della pipeline CI/CD; presentazione del progetto (demo live, poi un video narrato per la finale aziendale); \textbf{primo posto} a livello di ufficio, poi a livello aziendale su tutte le sedi Elca.
  \item \textbf{Sistema di Gestione Pensioni Svizzere:} Manutenzione e stabilizzazione di una suite di gestione pensionistica (backend \textbf{Spring}, frontend \textbf{JavaFX}) -- correzione bug e test backend con \textbf{Mockito}, con successiva estensione ad analisi dei task e supporto trasversale al team e ai business analyst.
  \item \textbf{Sistema di Prenotazione Rifugi Alpini (Club Alpino Svizzero \aux{\href{https://www.hut-reservation.org}{Link \linkicon}}):} Sviluppo e manutenzione delle funzionalità core di una piattaforma multilingua di prenotazione, utilizzata in oltre \textbf{500 rifugi alpini} (premio Gold \textbf{"Best Software Projects 2025"}, categoria Cloud Native Solutions) -- frontend \textbf{Angular}, backend \textbf{Spring}. Realizzazione di nuovi componenti, gestione dell'audit e generazione report Excel; sviluppo di suite di test \textbf{Jest} e \textbf{Cypress}, test di performance con \textbf{JMeter}, valutazione di una migrazione da Docker a Buildpacks, e migrazione della pipeline CI/CD da \textbf{Jenkins} a \textbf{GitHub Actions}.
\end{itemize}

\subsection{Progesi S.p.A. \rside{Roma}}
\subsubsection{Software Developer \rside{Ago 2023 -- Mar 2024}}
\begin{itemize}
  \item \textbf{Progetto Software Militare:} Sviluppo di un'applicazione desktop in \textbf{Java} per operazioni militari terrestri, in stretta collaborazione con il cliente principale, \textbf{Leonardo S.p.A.}; progettazione e implementazione di soluzioni software su misura per un contesto operativo altamente specialistico.
  \item \textbf{Installazione di Servizi Cloud Militari:} Configurazione e supporto di un servizio cloud avanzato per un centro militare, gestendo \textbf{web server, proxy e database} end-to-end, nel rispetto dei rigidi requisiti di sicurezza e affidabilità del settore.
\end{itemize}

\subsection{Join Business Management Consulting \rside{Roma}}
\subsubsection{Cyber Security Analyst \rside{Mag 2022 -- Ago 2023}}
\begin{itemize}
  \item \textbf{Analisi delle Applicazioni Bancarie:} Gestione autonoma di un'analisi completa di conformità degli accessi logici sul portafoglio applicativo di un importante gruppo bancario, valutando applicazioni interne ed esterne rispetto alla policy di sicurezza aziendale; sviluppo di tool \textbf{Python} personalizzati con interfaccia grafica per automatizzare l'analisi, con reporting dei risultati tramite presentazioni periodiche in Excel/PowerPoint al management.
  \item \textbf{Monitoraggio della Sicurezza:} Analisi quotidiana dei log \textbf{Netskope} per il rilevamento di accessi anomali, a supporto dei team \textbf{SCC} e \textbf{SOC} nella gestione degli incidenti di sicurezza.
\end{itemize}

\subsection{Solving Team \rside{Roma}}
\subsubsection{Corso Full Stack Developer \rside{Feb 2022 -- Mag 2022}}
\begin{itemize}
  \item Completamento di un percorso di formazione intensivo Full Stack Developer: realizzazione di un backend monolitico in \textbf{Spring} con funzionalità CRUD e gestione utenti, sviluppo di un frontend in \textbf{Angular}, ed esplorazione dei fondamentali \textbf{DevOps} incluso \textbf{Docker} -- basi tecniche su cui si fondano tutte le esperienze successive.
\end{itemize}

\section{Istruzione}
%==========
\subsection{Università degli Studi di Palermo}
\subsubsection{Laurea Triennale in Ingegneria Informatica \rside{2018 -- 2021}}
\vspace{1em}

\section{Competenze}
%===============
\subsection{Tecniche}
\begin{itemize}
  \item \textbf{Linguaggi:} Java, Python, TypeScript, JavaScript, SQL, Bash, C.
  \item \textbf{Framework e Librerie:} Spring Boot, Angular, FastAPI.
  \item \textbf{Dati e Messaggistica:} PostgreSQL, MySQL, MongoDB, Kafka.
  \item \textbf{Cloud e DevOps:} Azure, Docker, Git, GitHub Actions, Azure DevOps, CI/CD, amministrazione Linux.
  \item \textbf{Testing:} JUnit/Mockito, Jest, Cypress, JMeter.
  \item \textbf{AI:} integrazione di LLM in applicazioni e automazioni via API, prompt engineering, deployment di modelli in locale; utilizzo quotidiano di strumenti di coding assistito (Claude Code, GitHub Copilot).
  \item \textbf{Architetture e Metodologie:} microservizi, REST API, pipeline ETL, soluzioni cloud-native, lavoro in team Agile (daily, sprint planning).
\end{itemize}

\subsection{Lingue}
\begin{itemize}
  \item \textbf{Italiano}: Madrelingua.  
  \item \textbf{Inglese}: B1 (conoscenza intermedia).
\end{itemize}

\subsection{Interessi}
\noindent Sviluppo di \textbf{progetti personali} e \textbf{open-source}, con particolare interesse per l'\textbf{automazione} e il \textbf{self-hosting} (infrastruttura domestica autogestita). Speaker al \textbf{Career Day} dell'Università di Palermo (2025), con presentazione aziendale e demo tecnica dal vivo. Nel tempo libero, \textbf{numismatica} e \textbf{calcio}.




% \section{Certifications}
% %=======================
% \begin{enumerate}[label=\null, left=0pt..0pt, itemsep=0pt]
%   \item Some programming bootcamp, Location. \rside{2023}
%   \item Forbes Top 15 Procrastinators under 15. \rside{2022}
%   \item Perticipation award for attending workshop. \rside{2021}
% \end{enumerate}


% \section{Publications}
% %=====================
% \begin{enumerate}
%   \item A. Vaswani et al., “Attention is All you Need,” arXiv (Cornell University), vol. 30, pp. 5998 – 6008, Jun. 2017, [Online]. Available: \url{https://arxiv.org/pdf/1706.03762v5}
%   \item C. Darwin and J. Murray, \textit{On the origin of species by means of natural selection, or, The preservation of favoured races in the struggle for life.} 1859. doi: 10.5962/bhl.title.82303.
% \end{enumerate}


% \section{Awards}
% %===============
% \begin{description}
%   \item [2023] Some programming bootcamp, Location.
%   \item [2022] Forbes Top 15 Procrastinators under 15.
%   \item [2021] Conducted workshop on family planning by abstainance in SomePlace.
% \end{description}


% \section{Projects}
% %=================
% \subsection{Lunar Rover Project \aux{\href{https://robotics.nasa.gov/lmr/}{Link \linkicon}} \rside{2012 -- 13}}
% \begin{itemize}
%   \item LIDAR array, GPS, Variable FOV camera, Python, C/C++
%   \item Designed the rover to be modular, reconfigurable for sensor systems, handle science payloads, storage space, surplus resources, and robotics addons. % These modules shall be supported by centralized power and data ports. 
% \end{itemize}

% \subsection{Mars Rover \aux{\href{https://science.nasa.gov/mission/mars-exploration-rovers-spirit-and-opportunity/}{Project link \linkicon}} \rside{Feb 2012}}
% \begin{itemize}
%   \item Tech used, Language, Frameworks etc.
%   \item Implemented ABC feature
%   \item Optimised XYZ by 50\%
% \end{itemize}


\end{document}